\subsubsection{5.1 Anomalien}\label{anomalien}

Für die zusätzlich angenommene 3+1D-Weylinterpretation lautet das lokale
Anomaliepolynom bis auf Konventionsvorzeichen

\[\adjustbox{max width=\linewidth}{$\displaystyle 
 I_6=[\widehat A(T)\,\mathrm{ch}_R(F)]_6
 =\mathrm{ch}_3(F)-\frac{p_1(T)}{24}\,\mathrm{ch}_1(F).
$}\]

Der reine gravitative Anteil \([\widehat A]_6\) verschwindet; seine
Formgrade sind Vielfache von vier. Das ist die bekannte dimensionale
Unterscheidung bei
\href{https://collaborate.princeton.edu/en/publications/gravitational-anomalies/}{Álvarez-Gaumé
und Witten}. Eine Spur-/Weylanomalie ist nochmals ein anderer Begriff.

Auf (16,4) ist der SU(4)-Kubikkoeffizient 16 bei Normierung A(4)=1. Der
exakte Test \(t=\operatorname{diag}(1,1,1,-3)\) hat
\(\operatorname{tr}t=0\), \(\operatorname{tr}t^3=-24\). Ist SU(4)
\textbf{dynamisch geeicht}, muss die vollständige Theorie diese
Eichanomalie kompensieren. Als globale Symmetrie kann sie eine
't-Hooft-Anomalie tragen; dann folgt keine pauschale Spiegelpflicht. Ein
vollständiger konjugierter Spiegel ist ein möglicher Ausgleich, aber
hier nicht als einzige oder allgemein minimale Lösung bewiesen.

„Gravitativ 64`` kann sinnvoll eine \textbf{gemischte
U(1)-Gravitationsanomalie} meinen, wenn alle 64 linksgetragenen
Weylfelder explizit U(1)-Ladung eins erhalten. Das muss samt
Eich-/Globalstatus gesagt werden. Für SU(4) allein ist der gemischte
lineare Spurkoeffizient null. Die Quellenmarke N darf nicht ohne Beweis
in eine dynamisch geeichte chirale U(1) umgedeutet werden.

\subsubsection{5.2 Der behauptete Spin-2-Ausschluss ist
falsch}\label{der-behauptete-spin-2-ausschluss-ist-falsch}

Die Lorentzdarstellungen liefern exakt

\[\adjustbox{max width=\linewidth}{$\displaystyle 
 (\tfrac12,\tfrac12)\otimes(\tfrac12,\tfrac12)
 =(0,0)\oplus(1,0)\oplus(0,1)\oplus(1,1).
$}\]

Eine Ableitung des Vektorbilinears \(\psi^\dagger\bar\sigma_\mu\psi\)
kann daher einen symmetrischen spurfreien (1,1)-Tensor bilden.
Insbesondere hat der Energie-Impuls-Tensor
\(T_{\mu\nu}\sim i\psi^\dagger\bar\sigma_{(\mu}
\overleftrightarrow\partial_{\nu)}\psi\), mit passender Spurbehandlung,
Dimension \(3/2+3/2+1=4\). Ein Energie-Impuls-Tensor kann bereits in
einer festen Hintergrundraumzeit definiert werden; dynamische
Diffeomorphismen müssen nicht vorausgesetzt werden, um dieses
Gegenbeispiel zu formulieren.

Das schließt \textbf{T7 nicht}. Ein vorhandener Tensor ist kein
masseloser Spin-2-Pol. Das sinnvolle Ziel ist später sein transversaler,
spurfreier Korrelator auf demselben räumlichen Parent, mit positiver
Norm, beiden Helizitäten und universeller Kopplung. Wir entfernen hier
eine falsche Darstellungsobstruktion, nicht die dynamischen
Nachweispflichten.

\subsubsection{5.3 (1,0)-Kinetik ist nicht generell
verboten}\label{kinetik-ist-nicht-generell-verboten}

Der bilineare Ein-Ableitungs-Term mit einem (1,0)-Feld und seinem
Adjungierten besitzt keinen Lorentzskalar. Mit zwei Ableitungen
existiert aber etwa

\[\adjustbox{max width=\linewidth}{$\displaystyle 
 \mathcal L_2\propto
 (\partial^{\alpha\dot\alpha}\Phi_{\alpha\beta})
 (\partial^{\beta\dot\beta}\bar\Phi_{\dot\alpha\dot\beta})
$}\]

für symmetrisches \(\Phi\). Alle Spinorindizes sind kontrahiert. Bei
rein zeitartigem Impuls ist sein Symbol proportional zu
\(\omega^2(|\Phi_{11}|^2+2|\Phi_{12}|^2+|\Phi_{22}|^2)\), also nicht
identisch null. Das ist ein Gegenbeispiel zum uneingeschränkten Satz
„kein lokaler kovarianter kinetischer Term``. Es beweist \textbf{nicht}
Positivität des vollen Hamiltonoperators, korrekte Constraints,
gewünschte Freiheitsgrade oder eine native Quelle. Diese bleiben die
entscheidenden Tests.

\subsection{6. Die Vierzustandsrechnung richtig
lesen}\label{die-vierzustandsrechnung-richtig-lesen}

Das Modell mit \(s=\sqrt8g\), Bosonverbindung η und Diagonalen (0,Δ,Δ,0)
erlaubt kohärenten \textbf{Paartransport}. Die exakte Identität
\((H^3)_{41}=8g^2\eta\) stimmt. Unsere unabhängige Fließkomma-Abtastung
reproduziert auf \([0,3000]\) mit Schrittweite .05 den größten
gefundenen Wert .998920021869 bei 2484.05. Der erste
Gitter-Lokalhöchstwert liegt bereits bei 6.05. Weder „erster
Maximalzeitpunkt`` noch ein globales exaktes Maximum ist damit bewiesen.
Diese kleine Rechnung wird nicht mit dem vollständigen Einloch-Satz des
Hauptberichts verwechselt.

Für \(F_\pm(\epsilon)=(\epsilon\pm\sqrt{\epsilon^2+32g^2})/2\) gilt
\(F'_+-F'_-=\epsilon/\sqrt{\epsilon^2+32g^2}\). Bei ε=0 stimmen die
Ableitungen überein, aber die \textbf{Eigenwertlücke} beträgt
\(\sqrt{32}|g|>0\). Ein verschwindender nackter Vermittlerparameter und
eine verschwindende wechselwirkende Anregungslücke sind verschieden. Ein
allgemeiner relativistischer Kegelsatz oder ein zwingendes neues
Renormierungsprogramm folgt aus dieser Ableitungsidentität nicht.

\subsection{7. Konsequenz für die nächste
Forschungsrevision}\label{konsequenz-fuxfcr-die-nuxe4chste-forschungsrevision}

Die sinnvollen positiven Anschlüsse sind jetzt schärfer:

\begin{enumerate}
\def\labelenumi{\arabic{enumi}.}
\tightlist
\item
  Den expliziten Ordnung-vier-Lift mit dem tatsächlich markierten Clock
  und der Herkunft des A3-Index vergleichen: vollständige Intertwiner,
  Ladungsphase und Bosonwirkung, nicht nur Gruppennamen.
\item
  Einen Quellenoperator finden, der die lokale Parität beider
  operational bestimmten Teile gleichzeitig ändert. Die native
  Z4-Symmetrie selbst tut das nicht. Der kontrollierte Einloch-Transfer
  ist bereits als Akzeptanztest für einen solchen Operator verfügbar.
\item
  Den Zwei-Ableitungs-Feldkandidaten auf Positivität und Constraints
  prüfen, bevor er als relativistischer Adapter zählt. Ein bloßes
  Nichtnull-Symbol ist noch keine gesunde Feldtheorie.
\item
  Den Energie-Impuls-Tensor als zulässigen \textbf{Kandidaten} behalten,
  aber erst auf dem gemeinsam konstruierten räumlichen Träger nach einem
  dynamischen Spin-2-Pol suchen. Kein vorhandener Spin-2-No-go
  rechtfertigt hier das Aufgeben dieses Anschlusses.
\end{enumerate}

62 zusätzliche Prüfbedingungen bestehen normal und optimiert identisch.
Exakte Identitäten, bedingte Darstellungsargumente und die numerische
Zeitfenstersuche sind im Ergebnisbericht getrennt. Der externe Status
„alle Restfragen geschlossen oder auf genau ein Stück reduziert`` wird
nicht übernommen. Die vollständige TOE bleibt offen.

\begin{center}\rule{0.5\linewidth}{0.5pt}\end{center}

\section{Historischer Beweisanhang: vollständiger nativer
Herleitungsstand
v1.6.4}\label{historischer-beweisanhang-vollstuxe4ndiger-nativer-herleitungsstand-v1.6.4}

Dieser Anhang bewahrt die frühere Herleitung vollständig. Seine
damaligen Statussätze werden durch die aktuelle Konsolidierung und den
Nachtrag oben ergänzt; er ist kein ungeprüft übernommener neuer
Gesamtstatus. Die externe laufende Symmetrieergänzung wurde separat
geprüft und ist nicht unbemerkt in diesen eingefrorenen Anhang
eingegangen.

\section{TFPT / Universalraum: nativer Pol, Bewegungsgleichung und
minimale
Feldtypen}\label{tfpt-universalraum-nativer-pol-bewegungsgleichung-und-minimale-feldtypen}

\textbf{Konsolidierte Forschungsfortsetzung v1.6.4 · 15. September 2026}

Diese Revision verbindet den während der Arbeit neu eingegangenen
Polsatz v1.6.3 mit einer unabhängig begonnenen Untersuchung der nativen
Bewegungsgleichung. Sie ergänzt die bisherigen Hauptdokumente; sie ist
keine verkürzte Neufassung des vollständigen Hauptbuchs. Die vorhandenen
Haupt- und Update-PDFs wurden in dieser Runde nicht verändert.

\subsection{1. Ergebnis in einem Satz}\label{ergebnis-in-einem-satz}

Auf dem abgesicherten Grundzustand des festgelegten nativen Fockmodells
existiert eine isolierte ursprüngliche Fermion-Entnahmelinie; nach
Kombination beider Rechnungen trägt sie \textbf{mehr als 88,007628 \%
des gesamten normierten Spektralgewichts pro Mode}. Die zusätzliche
Antwort ist jedoch nicht exakt auf eine einzige weitere Linie
reduzierbar. Ihr erster Rückwirkungsschritt wird direkt von derselben
ursprünglichen Wechselwirkung bestimmt.

Am Prüfpunkt \(g/\Delta=1/20\), ausdrücklich ohne \(\mu N\)-Zusatz:

{\def\LTcaptype{none} % do not increment counter
\begin{longtable}[]{@{}
  >{\raggedright\arraybackslash}p{(\linewidth - 4\tabcolsep) * \real{0.2727}}
  >{\raggedleft\arraybackslash}p{(\linewidth - 4\tabcolsep) * \real{0.3636}}
  >{\raggedleft\arraybackslash}p{(\linewidth - 4\tabcolsep) * \real{0.3636}}@{}}
\toprule\noalign{}
\begin{minipage}[b]{\linewidth}\raggedright
Größe
\end{minipage} & \begin{minipage}[b]{\linewidth}\raggedleft
Eingegangene v1.6.3
\end{minipage} & \begin{minipage}[b]{\linewidth}\raggedleft
Konsolidierte strengere Grenze
\end{minipage} \\
\midrule\noalign{}
\endhead
\bottomrule\noalign{}
\endlastfoot
Mittlere Bosonenzahl & \(0.77<\bar b<1.45\) &
\(0.842846<\bar b<1.245656\) \\
Energie der niedrigen Entnahmelinie &
\(0.007737<\epsilon/\Delta<0.062277\) &
\(0.007737<\epsilon/\Delta<0.039079764\) \\
Gewicht dieser Linie im gesamten CAR-Maß &
\(Z_{\rm low}>0.864972353\ldots\) &
\(Z_{\rm low}>0.880076280689\ldots\) \\
Obere Gewichtsgrenze & \(Z_{\rm low}<0.9759375\) &
\(Z_{\rm low}<0.9736610625\) \\
Übrige Entnahmeenergien & \(>0.379636\Delta\) & unverändert \\
Sämtliche Additionsenergien & \(>0.329636\Delta\) & unverändert \\
\end{longtable}
}

Die Dezimalzahlen sind gerundete Darstellungen rationaler Schranken. Es
sind weder exakte Polpositionen noch angepasste Zentralwerte. Der Pol
steht im retardierten Spektrum bei negativer Frequenz \(-\epsilon\). Das
niedrige Niveau im N=63-Hilbertraum ist 64-fach entartet; jede diagonale
Modenantwort sieht dieselbe Linie.

\textbf{Nicht bewiesen:} die physische Herleitung des Hamiltonoperators,
sein vollständiger Operationssatz, native Präparation, räumliche
Ausbreitung, ein vollständiges relativistisches Feldwörterbuch oder eine
TOE. T1-T8 bleiben als vollständige Aufgaben offen.

\subsection{2. Ein unveränderter
Modellvertrag}\label{ein-unveruxe4nderter-modellvertrag}

Es werden keine Hopping-, Massen-, Ladungs- oder Projektionsglieder zu H
hinzugefügt:

\[\adjustbox{max width=\linewidth}{$\displaystyle 
H=\Delta N_b+g(Q_++Q_-),\quad Q_+=\sum_A b_A^\dagger P_A,
\quad Q_-=Q_+^\dagger,
$}\] \[\adjustbox{max width=\linewidth}{$\displaystyle 
P_A=\sum_{i<j}W_{A,ij}f_jf_i,\quad N=N_f+2N_b,
\quad WW^\dagger=8I_{60}.
$}\]

Es gibt 64 CAR-Fermionmoden, 60 CCR-Bosonmoden und 480 von null
verschiedene reelle W-Einträge mit ihren ursprünglichen Vorzeichen.
\(\Delta>0\), g ist reell; die Zahlen beziehen sich auf
\(g/\Delta=1/20\). Die Tensorquelldatei ist durch SHA-256
\codeword{3f00a0892157a691e4ef26920c3702772c7bf76a250fad98ff04a079bb64b763}
fixiert.

Die Fockrealisierung, H, der sektorenübergreifende Energievergleich und
der Prüfparameter sind weiterhin Modellvoraussetzungen. Die geometrische
oder arithmetische Herkunft eines W-Tensors leitet diese Voraussetzungen
nicht allein her. Insbesondere sind die fünf elementaren CAR in einer
Spinor-Matrixkonstruktion nicht mit den hier verwendeten 64 Fockmoden
gleichzusetzen.

\subsection{3. Was tatsächlich erneut geprüft
wurde}\label{was-tatsuxe4chlich-erneut-gepruxfcft-wurde}

\subsubsection{3.1 Grundzustand: frische Berechnung statt bloßer
Statusübernahme}\label{grundzustand-frische-berechnung-statt-blouxdfer-statusuxfcbernahme}

Die ursprünglichen Programme wurden mit fixierten Dateihashes in eine
eigene Arbeitskopie übernommen. Die Paar-Konfigurationen wurden frisch
mit neu kompiliertem Quellcode berechnet:

{\def\LTcaptype{none} % do not increment counter
\begin{longtable}[]{@{}
  >{\raggedright\arraybackslash}p{(\linewidth - 4\tabcolsep) * \real{0.2727}}
  >{\raggedleft\arraybackslash}p{(\linewidth - 4\tabcolsep) * \real{0.3636}}
  >{\raggedleft\arraybackslash}p{(\linewidth - 4\tabcolsep) * \real{0.3636}}@{}}
\toprule\noalign{}
\begin{minipage}[b]{\linewidth}\raggedright
Ordnung
\end{minipage} & \begin{minipage}[b]{\linewidth}\raggedleft
Vollständig enumerierte Bosonkonfigurationen
\end{minipage} & \begin{minipage}[b]{\linewidth}\raggedleft
Normquadrat von \(Q_+^kF\)
\end{minipage} \\
\midrule\noalign{}
\endhead
\bottomrule\noalign{}
\endlastfoot
0 & analytischer Ausgangszustand & 1 \\
1 & 60 Kanäle / 480 Paare & 480 \\
2 & 1830 & 439680 \\
3 & 37820 & 575078400 \\
4 & 595665 & 952296652800 \\
\end{longtable}
}

Zusätzlich wurden unabhängige Wortentwicklungen, Überlaufgrenzen, die
vollständige Zweikörper-Casimiridentität, rationale Sektorvergleiche und
die unabhängige Symmetriespurrechnung wiederholt. Die vier ausgewählten
ursprünglichen Zertifikate bestehen normal und unter \codeword{-OO} mit
identischen JSON-Bytes. Dies ist eine gezielte erneute
Grundzustandsprüfung, nicht die Behauptung, alle historischen
Repository-Tests erneut ausgeführt zu haben. Der unveränderte alte
Normprüfer gibt unter dem aktuellen NumPy eine ComplexWarning bei der
Ganzzahlkonversion aus. Die neuen Prüfer bestätigen vor jeder Konversion
exakt verschwindende Imaginärteile und ganzzahlige Realteile des
gepinnten W. Die Warnung ist gespeichert und wird nicht als verlorener
Tensoranteil oder als unterdrückter Prüffehler ausgegeben.

Der modellinterne Satz bleibt bestehen: Für \(0<|g|/\Delta\le1/20\) ist
der globale Grundzustand \(\Omega\) eindeutig, hat N=64 und ist
Spin(10)×SU(4)-invariant. Am Zahlenprüfpunkt:

\[\adjustbox{max width=\linewidth}{$\displaystyle 
-1.158089\Delta<E_0<-1.129636\Delta,
\qquad \operatorname{gap}(H)>0.007737\Delta.
$}\]

Der volle Grundzustandsvektor ist damit nicht ausgerechnet. Die fünf
Versuchsvektoren sind keine behauptete invariante Fünferbasis für H.

\subsubsection{3.2 Das neu eingegangene
Polergebnis}\label{das-neu-eingegangene-polergebnis}

Der vollständige technische Bericht v1.6.3 und sein 248-zeiliger Prüfer
wurden gelesen; Quellen, Bericht, Ergebnismatrix und Prüfpaket wurden
unverändert archiviert. Der Prüfer wurde normal und optimiert
wiederholt: \textbf{317 Prüfbedingungen, identische JSON-Bytes auch zum
gelieferten Bericht}.

Der Herkunftspin des Prüfers ist
\codeword{fdbcabd244450c302182086d67c68284634c994e3fee026200c42c508f731654}.
Er verwendet denselben ursprünglichen
\codeword{verify\_hole.py}-Quellstand wie die bisherige Fortsetzung. Die
neuen eigenen Sektorvergleiche wurden zusätzlich unabhängig mit
rationalen LDL-Pivots berechnet.

Die Prüfzahlen zählen auch Komponenten und Wiederholungen; sie sind
keine Zahl unabhängiger Entdeckungen. Die analytischen Argumente sind
ausgeschrieben, aber nicht vollständig in einem Beweisassistenten
formalisiert.

\subsection{4. Der einfache native Anschluss: Bewegungsgleichung statt
Umdeutung}\label{der-einfache-native-anschluss-bewegungsgleichung-statt-umdeutung}

Erweitere jede W-Zeile zur antisymmetrischen Matrix \(M_A\) mit
\((M_A)_{ij}=W_{A,ij}\) für i\textless j. Definiere

\[\adjustbox{max width=\linewidth}{$\displaystyle 
D_r=\sum_{A,j}(M_A)_{rj}b_Af_j^\dagger,
\qquad \phi_r=D_r/\sqrt{15}.
$}\]

Direkte CAR/CCR-Normalordnung liefert \textbf{Operatoridentitäten auf
dem endlichen Teilchenkern}, nicht nur Gleichheiten auf Testzuständen:

\[\adjustbox{max width=\linewidth}{$\displaystyle 
\boxed{[H,f_r]=-gD_r,\qquad [H,f_r^\dagger]=gD_r^\dagger,}
$}\] \[\adjustbox{max width=\linewidth}{$\displaystyle 
[N,D_r]=-D_r,\qquad \{f_r,D_s^\dagger\}=0.
$}\]

Der neue Vergleichsoperator hat also dieselbe Ladung −1 wie f. Er ist
nicht das früher auf dem leeren Referenzzustand betrachtete
\(\chi^\dagger\sim b^\dagger f^\dagger\) mit Ladung +3. Auf dem
Grundzustand liegen \(f\Omega\) und \(D\Omega\) in N=63,
\(\chi^\dagger\Omega\) dagegen in N=67. Neutrale Zeitentwicklung hebt
diese Unterscheidung nicht auf.

Die Normalordnungsprüfung erfasst alle 64 ursprünglichen f-Operatoren
und die tatsächlichen W-Vorzeichen. Für die Kompositnorm gilt exakt

\[\adjustbox{max width=\linewidth}{$\displaystyle 
\{D_r,D_s^\dagger\}=
\sum_{A,B,j,k}(M_A)_{rj}(M_B)_{sk}
\left(\delta_{AB}f_j^\dagger f_k+\delta_{jk}b_B^\dagger b_A\right).
$}\]

Mit \(\bar b=\langle N_b\rangle\) und der Grundzustandssymmetrie:

\[\adjustbox{max width=\linewidth}{$\displaystyle 
\langle\{D_r,D_s^\dagger\}\rangle=\delta_{rs}S,
\qquad S=15-\frac7{32}\bar b.
$}\]

\(\phi\) hat entsprechend Norm \(1-7\bar b/480\), nicht globale
kanonische CAR. Der normierte Zustandserwartungswert ersetzt keine
Operatorrelation.

